\documentclass[11pt]{article}
\usepackage[T1]{fontenc}
\usepackage{lmodern}
\usepackage{royden-report}
\usepackage{array}
\title{Long-distance water-wave Dysthe experiments}
\author{Experimental protocol}
\date{26 October 2026}
\newcommand{\norm}[1]{\left\lVert #1\right\rVert}
\newcommand{\Dabs}{\lvert D_\tau\rvert}
\begin{document}
\maketitle

\noindent\textbf{Status.} Proposed cluster experiments, following a completed CPU
pilot. This document contains no A100 training results and no claim of global
existence. Its objective is to determine how far resolved numerical solutions
can be followed, and how accurately a slab PINN reproduces their nonlinear
evolution across several initial wave profiles. A failed accuracy or resolution
check is a result to retain, not a reason to omit a case.

\section{Equation and meaning of distance}

The approved model, \texttt{water-wave-dysthe-spatial-v1}, is the spatial
free-surface-envelope equation
\begin{equation}\label{eq:model}
\begin{split}
u_\xi={}&-\mathrm{i}u_{\tau\tau}-\mathrm{i}|u|^2u\\
&-\epsilon\bigl(8|u|^2u_\tau+2u^2\overline{u_\tau}
                  +2\mathrm{i}u\Dabs(|u|^2)\bigr).
\end{split}
\end{equation}
Here \(u\colon[0,\Xi]\times\mathbb R/(L\mathbb Z)\to\mathbb C\),
\(\xi\) is propagation distance, and \(\tau\) is retarded time. The Fourier
multiplier of \(\Dabs\) is \(|k|\), with zero value at \(k=0\).
The free-surface coefficient of \(u^2\overline{u_\tau}\) is \(2\); a
potential-envelope equation is a different model. Equation~\eqref{eq:model}
and its formal weakly nonlinear, narrowband setting follow the spatial
formulation of Fedele and Dutykh~\cite{fedele}.

If \(X\) and \(T\) denote dimensional distance and time, the carrier
coordinates satisfy \(x=k_0X\), \(t=\omega_0T\). The present normalization is
\[
\xi=\epsilon^2k_0X,\qquad
\tau=\epsilon(2k_0X-\omega_0T),\qquad B=\epsilon u.
\]
Consequently, reporting metres requires a specified carrier wavenumber
\(k_0\) and \(\epsilon>0\). The limit \(\epsilon=0\) in
\eqref{eq:model} defines an NLS control in the same computational coordinates;
it is not an invertible physical coordinate transformation at zero steepness.

A constant profile of amplitude \(A>0\) has solution
\(u(\xi)=A\exp(-\mathrm{i}A^2\xi)\). Define the nonlinear distance scale by
\[
L_{\mathrm{NL}}=A^{-2},\qquad n_{\mathrm{NL}}(\xi)=A^2\xi.
\]
One such length corresponds to one radian of uniform-wave nonlinear phase;
it is not a recurrence period. The machine diagnostics use
\(A=\norm{u(0)}_\infty\), explicitly labelled as a peak-based comparison scale.
For sideband perturbations this differs from the prescribed background
amplitude; background-based distances must be labelled separately if reported.
For example, if \(A=0.5\), the distances \(4,8,16,32,64\) correspond to
\(1,2,4,8,16\) nonlinear lengths. These are research horizons, not guaranteed
physical validity intervals.

\section{Cases and staged comparisons}

Use the new A100 campaign's 36 deterministic initial-condition groups:
12 localized packets, 12 sideband-perturbed waves, and 12 random band-limited
profiles. The existing split assigns 24 groups to training, six to validation,
and six to testing. The long-distance study evaluates all 12 validation and
test groups; refinements and perturbations retain their parent group identity.
For PINNs, these names define evaluation cohorts: every new inlet is fitted
individually, without reference trajectory labels. They do not imply learned
operator generalization.

\begin{center}
\begin{tabular}{@{}>{\raggedright\arraybackslash}p{0.23\linewidth}>{\raggedright\arraybackslash}p{0.20\linewidth}>{\raggedright\arraybackslash}p{0.48\linewidth}@{}}
\toprule
Experiment & Horizons \(\Xi\) & Question answered\\
\midrule
Reference qualification & \(4,8,16,32,64\) & Can the same trajectory be
resolved independently in profile grid and propagation step?\\[3pt]
Slab PINN & \(4,8,16\) & How far does a neural approximation remain accurate
without resets to reference data?\\[3pt]
NLS control & Matched reference horizons & Which changes arise from the
higher-order Dysthe terms?\\[3pt]
Sensitivity & \(4,8,16\) & Do nearby inlets separate
reproducibly under a specified perturbation?\\[3pt]
Packet domain check & \(4,8,16,32,64\) & Are periodic copies
affecting the localized-packet evolution?\\
\bottomrule
\end{tabular}
\end{center}

Numerical and neural horizons must be reported separately. A resolved
reference at \(\xi=64\) does not validate a PINN beyond its evaluated interval.
Each PINN horizon is fitted independently from the same original inlet;
checkpoint restart continues that fit, while the next horizon starts its own
declared fit. There are 108 proposed fits: 12 inlets, three seeds and three
horizons. A longer fit receives proportionally more total updates.
Longer horizons are interpreted only while their independent checks pass.
Preserve the first failure distance and the last qualified horizon. Retain
all seeds, including failed fits; do not extend only successful cases and then
report the surviving subset as the original campaign.

\section{Numerical reference and diagnostics}

\subsection{Resolution and regularity}

Use the validated Fourier reference solver with integrating-factor RK4.
For each case, use \(N=1024\), step \(h=0.00125\), and stored spacing
\(\Delta\xi=0.05\). Compare independently against \(N/2\) at the same step
and against \(2h\) at the same grid. Compare on common stations. Let \(v\)
be the refined reference and \(w\) a
comparison. For \(s=0,1,2\), record
\[
E_s(\xi)=\frac{\norm{w(\xi)-v(\xi)}_{H^s}}
                       {\norm{v(\xi)}_{H^s}},\qquad
\norm{v}_{H^s}^2=L\sum_k(1+k^2)^s|\widehat v_k|^2.
\]
Here Fourier coefficients use the \(1/N\) forward normalization. Decisions
use the maximum over evaluated stations, not only the endpoint. Each independent
comparison must satisfy \(E_0\leq10^{-5}\), \(E_1\leq10^{-4}\) and
\(E_2\leq10^{-3}\). Also require action drift at most \(10^{-6}\) and spectral
tail action fraction at most \(10^{-8}\), where the tail consists of modes
at or above 80\% of the declared Fourier cutoff. These are numerical study tolerances,
not theorem hypotheses or physical thresholds.

Track \(\norm{u}_\infty\), \(\norm{u}_{H^1}\),
\(\norm{u}_{H^2}\), and
\(R_s(\xi)=\norm{u(\xi)}_{H^s}/\norm{u(0)}_{H^s}\).
Record spectral-edge energy and mass drift as complementary diagnostics.
A finite-dimensional Fourier solution is smooth by construction; bounded
sampled norms on one grid do not certify continuum regularity. Increasing
derivative norms require refinement and tail inspection before physical
interpretation. Conservation alone does not establish trajectory accuracy.

\subsection{Focusing, recurrence and spectrum}

Record the maximum amplitude against distance, candidate focusing locations,
spectral mean, spectral width, sideband asymmetry, and distances from the
initial profile. Distinguish raw complex-field error from a diagnostic that
removes a global phase; alignment must never conceal
phase errors in the main accuracy gate. If repeated focusing and recovery
occur, determine whether event locations and amplitudes survive refinement.
The target is to resolve repeated events where they exist, not to manufacture
a recurrence count in every case. Armaroli, Brunetti and Kasparian~\cite{armaroli}
motivate examining recurrence and transient spectral shifts in deep-water
Dysthe dynamics; their reduced-model results do not predict a universal
recurrence interval for these inlets.

The event finder marks interior amplitude peaks with prominence at least 2\%
of the initial peak. Recurrence candidates are local minima of phase-aligned
relative \(L^2\) distance no larger than 5\%, with prominence 0.02. These
thresholded, sampled events are candidates for refined inspection, not exact
continuous event times; no translation is removed.

Run the NLS control with identical inlet, computational coordinates and
numerical qualification. Compare the full complex field, focusing events and
spectra. A discrepancy measures the consequence of the model terms at these
parameters; it does not by itself establish that Dysthe is closer to a
laboratory wave or the full water-wave equations.

\subsection{Sensitivity and packet-domain checks}

For a recorded smooth direction \(h\), normalized by
\(\norm{h}_{H^1}=1\), compute the centered directional response
\[
q_\delta(\xi)=\frac{S_\xi(u_0+\delta h)-S_\xi(u_0-\delta h)}{2\delta},
\qquad G_\delta(\xi)=\norm{q_\delta(\xi)}_{H^1}.
\]
Use \(\delta=10^{-3}\). Compare \(\delta\) and \(\delta/2\), then independently refine the numerical
resolution for the signed response \(q_\delta\). Agreement of amplification
norms alone can hide an incorrect response direction. The result concerns one
declared direction per case, not the operator norm of a solution-map derivative
and not a Lyapunov exponent. Growth can be genuine finite-distance instability.
The signed-response comparison must agree within 1\% in \(H^1\) under both
the perturbation-size and numerical-resolution checks.

For localized packets, double the periodic domain while retaining physical
profile spacing and the same localized inlet definition. Compare over the
common central region and inspect boundary energy. A large discrepancy means
the smaller periodic domain is inadequate for a whole-line interpretation;
it does not automatically make its periodic solution incorrect. Sideband and
random periodic waves are not silently reinterpreted as localized packets.
Require initial central-profile consistency within 0.1\% in \(H^1\), then
evolution agreement within 1\% in \(H^1\). Both domains must independently
pass their grid/step qualifications before the comparison is interpreted.

\section{Slab PINN experiment}

Divide propagation into slabs of width \(0.5\). Train one spectral PINN on
each slab, using its predicted endpoint as the next inlet. Every slab retains
the complete periodic \(\tau\)-profile for the nonlocal term. The model uses
a hard inlet constraint, exact linear propagation and a learned Fourier
correction; a small neural network maps local distance to complex Fourier
coefficients. Only the first inlet is prescribed from the case data. No slab
interface is reset to the numerical trajectory.

The proposed A100 setting fixes width 64, depth 3, output modes
\(|m|\leq48\), 256 profile points, 64 distance collocation slices per update,
and 400 Adam updates per slab. The intervals \(\Xi=4,8,16\) therefore contain
8, 16 and 32 slabs and allow 3,200, 6,400 and 12,800 updates respectively.
Use optimizer seeds 11, 29 and 47 on each evaluation inlet. This fixed work
per unit distance makes longer runs interpretable; it does not promise that
the training budget remains adequate when waves become more complicated.

Training and checkpoint selection use the full Dysthe residual, without
reference trajectory labels. Evaluate residuals on an independent denser
set of stations and profile points. Residual size is diagnostic and is not
substituted for solution error. At every stored station, require
\[
E_0\leq0.01,\quad E_1\leq0.05,\quad E_2\leq0.05,\quad
\frac{|\norm{u_{\mathrm{PINN}}}_\infty-
         \norm{u_{\mathrm{ref}}}_\infty|}
        {\norm{u_{\mathrm{ref}}}_\infty}\leq0.05.
\]
The \(H^2\) condition is an additional long-distance criterion; the completed
CPU pilot used \(L^2\), \(H^1\) and peak-amplitude gates. Report all four
quantities separately for every case and seed, together with
the worst-case and median summaries, first failed station, runtime, and peak
GPU memory. A failed reference qualification invalidates a numerical accuracy
claim; it is not counted as a PINN success or failure against trustworthy truth.

The CPU pilot motivated this neural candidate, but the conventional solver
remained faster and more accurate. Measure complete fitting cost for every
new inlet, not only evaluation of an already fitted network. A GPU comparison
must state hardware, precision, synchronization and whether setup/data
generation are included. FNO comparisons from the earlier small-data pilot
remain separate; this experiment is not a new operator-learning benchmark.

\section{Decision and interpretation}

The principal outcome is a table of the last qualified numerical horizon and
last accurate PINN horizon for each inlet and seed, expressed both in \(\xi\)
and in the stated nonlinear-length scale. Secondary outcomes are refined
focusing and spectral diagnostics, directional sensitivity, NLS differences,
and packet-domain sensitivity. Report diagnostic availability and failed checks
explicitly. No predetermined bound such as \(\norm{u}_{H^2}<5\) is an objective.

Monitor \(\epsilon\norm{u}_\infty\) and spectral bandwidth to expose movement
outside the weakly nonlinear, narrowband regime. These are model-validity
indicators; no universal numerical cutoff makes the asymptotic approximation
valid for arbitrary distances. The experiments address a periodic reduced
water-wave model, not wave breaking, dissipation, transverse dynamics or
full Euler evolution.

\begin{remark}
Resolved finite-distance evolution and accurate neural approximation are
empirical evidence. They do not prove global existence, exclude later
singularity formation, or establish convergence as resolution tends to
infinity. A rigorous existence certificate would require an applicable
analytical or validated-numerical argument with controlled residual,
stability and unresolved-tail bounds.
\end{remark}

The runnable configuration \texttt{configs/long-distance-a100.json} and cluster
guide are the authoritative execution specification. A separate smoke
configuration exercises the workflow and does not establish long-distance
accuracy. The command module \texttt{dysthe\_experiments.long\_distance} has
\texttt{reference}, \texttt{pinn}, \texttt{sensitivity}, \texttt{domain} and
\texttt{summarize} stages. The cluster submission helper defaults to a dry run;
actual submission requires its explicit \texttt{--submit} option. No cluster
job has been run to produce this protocol.

Preserve configuration hashes, code revision, environment, inlet hashes,
checkpoints and raw traces. Publish compact summaries and failure reasons;
keep large arrays and checkpoints outside Git. This protocol is editable
LaTeX and is distributed with its style file and rendered PDF.

\begin{thebibliography}{9}
\raggedright
\bibitem{fedele} F.~Fedele and D.~Dutykh, \emph{Hamiltonian description and
traveling waves of the spatial Dysthe equations}, research report (2011;
revised 2012), \href{https://arxiv.org/abs/1110.3605}{arXiv:1110.3605}.
\bibitem{armaroli} A.~Armaroli, M.~Brunetti and J.~Kasparian,
\emph{Recurrence in the high-order nonlinear Schr\"odinger equation:
A low-dimensional analysis}, Physical Review E \textbf{96}, 012222 (2017),
\href{https://doi.org/10.1103/PhysRevE.96.012222}{doi:10.1103/PhysRevE.96.012222}.
\end{thebibliography}
\end{document}
